\subsection{Simplifying Expressions} 
\begin{definition}[Simplify]
To \term{simplify} an expression means to manipulate it to get an \term{equivalent expression} in some standard form. Expressions are \term{equivalent} if they represent the same number, \emph{no matter what} the values of the variables.
\end{definition}
To learn to simplify a new kind of expression, you must learn:
\begin{itemize}
\item How to manipulate an expression of that kind.
\item What the standard form is.
\end{itemize}
We already know how to manipulate expressions in which variables are only added, subtracted, and multiplied by real numbers, which we call \makeblank[1in] expressions.
\begin{Example}
Simplify $3(2x-5)+4x-11z+5+5(z-1)$ as much as possible. Justify each step with one of the properties of real numbers.
\begin{cpic}{}{4}
\end{cpic}
To check that these expressions are equivalent, we can pick any two real numbers for $x$ and $z$. Let's choose $x=\makeblanksmall$ and $z=\makeblanksmall$.
\begin{lpic}{}{8}
\regtextleft{0}{-1}{$3(2x-5)+4x-11z+5+5(z-1)=$};
\regtextleft{0}{-6}{$10x-6z-15=$};
\end{lpic}
\end{Example}